\documentclass[11pt]{amsart}
\usepackage{amsmath,amssymb,amsthm,mathtools}
\usepackage[margin=1in]{geometry}
\usepackage{hyperref}

\newtheorem{theorem}{Theorem}[section]
\newtheorem{proposition}[theorem]{Proposition}
\newtheorem{lemma}[theorem]{Lemma}
\newtheorem{corollary}[theorem]{Corollary}
\theoremstyle{definition}
\newtheorem{definition}[theorem]{Definition}
\theoremstyle{remark}
\newtheorem{remark}[theorem]{Remark}

\newcommand{\RH}{\mathrm{RH}}
\renewcommand{\Re}{\operatorname{Re}}
\renewcommand{\Im}{\operatorname{Im}}
\newcommand{\half}{\tfrac12}
\newcommand{\fhat}{\widehat f}
\newcommand{\Zoff}{\mathcal Z_{\mathrm{off}}}
\newcommand{\Zcl}{\mathcal Z_{\mathrm{cl}}}
\newcommand{\xio}{\xi_0}
\newcommand{\Pmpoly}{P_{4m}}
\newcommand{\psiO}{\psi_0}
\newcommand{\cN}{\mathcal N}
\newcommand{\kap}{\kappa}
\DeclareMathOperator{\sgn}{sgn}

\title[The Defect of Off-Line Zeros]{The arithmetic defect of off-line zeros:\\
parity, the modified explicit formula, and two-scale structure}
\author{David Alejandro Trejo Pizzo}
\thanks{Independent Researcher. \texttt{dtrejopizzo@gmail.com}}


\begin{document}
\begin{abstract}
We study the exact arithmetic footprint that a \emph{finite} set of off-critical-line zeros
of the Riemann zeta function would leave. The off-line zeros of $\zeta$ form complete Klein
four-orbits $\{\half\pm b\pm i\gamma\}$, and the negative index $\kappa$ of the localized
Weil form is the dimension of its maximal negative subspace. We first prove the geometry of
that subspace: each orbit contributes exactly two linearly independent negative directions,
so $\kappa=2m$ is even, with the minimal non-trivial value $2$; directions from orbits at
distinct heights are exponentially nearly orthogonal, so the negative subspace decomposes as
an approximate orthogonal sum of Lorentz planes. Assuming a finite negative index
(Hypothesis~D, $0<\kappa=2m<\infty$, strictly weaker than the negation of the Riemann
Hypothesis), the completed zeta factors as $\xi=\xi_0\,P_{4m}$ with $P_{4m}$ the degree-$4m$
polynomial of the off-line zeros and $\xi_0$ an entire function sharing the functional
equation and having all zeros on the line. This yields an \emph{exact modified explicit
formula} $\psi(x)=\psi_0(x)-D_m(x)$, where the arithmetic defect $D_m$ is a finite sum of
$4m$ oscillations whose amplitudes, frequencies, and phases are explicit; its dominant term
grows like $x^{1/2+b_{\max}}$ and its energy like $X^{2b_{\max}}$. We then locate the precise
obstruction: the companion $\zeta_0=\xi_0/(\pi^{-s/2}\Gamma(s/2))$ has poles inside the
critical strip, so by Landau's theorem it admits no Dirichlet series with non-negative
coefficients and no Euler product --- the analytic factorization is incompatible with the
multiplicative structure of $\zeta$. This is the wall we name (RSRP): the Euler product is
incompatible with a finite non-zero index; it follows from the Riemann Hypothesis, does not
imply it, and is distinct from the positivity, de Branges, and pair-correlation walls.
Finally we exhibit the two-scale signature of a single orbit --- an exponentially growing
oscillation in the Chebyshev variable and a localized logarithmic depression of
$\log|\zeta(\half+it)|$ near $t=\gamma$ --- prove that both are compatible with the
unconditional prime number theorem and (in the polynomially-bounded regime) with the
multiplicative-chaos model, and state precisely which three inputs any proof of the exclusion
principle $\kappa\in\{0,\infty\}$ would require, none of which is currently available.
Nothing here assumes or proves the Riemann Hypothesis.
\end{abstract}
\maketitle

\tableofcontents
\newpage

\section{Introduction}

The Riemann Hypothesis ($\RH$) is the assertion that every nontrivial zero
$\rho=\beta+i\gamma$ of $\zeta$ has $\beta=\half$. Suppose it fails. What, precisely, would a
finite collection of off-line zeros do to the distribution of primes, and where exactly does
the assumption that they exist collide with the arithmetic of $\zeta$? This paper answers both
questions exactly, under the working hypothesis that the index is finite.

We write $\xi(s)=\half s(s-1)\pi^{-s/2}\Gamma(s/2)\zeta(s)$ for the completed zeta function,
$\Zoff=\{\rho:\beta_\rho\ne\half\}$ for the off-line zeros and $\Zcl=\{\rho:\beta_\rho=\half\}$
for the on-line ones. The functional equation $\xi(s)=\xi(1-s)$ and the reality
$\overline{\xi(\bar s)}=\xi(s)$ make $\Zoff$ invariant under the Klein four-group
$V=\{1,\sigma,\iota,j\}$, $\sigma(s)=\bar s$, $\iota(s)=1-s$, $j=\sigma\iota$, so the off-line
zeros come in quartets
\[
\mathcal O_j=\{\sigma_j\pm i\gamma_j,\ (1-\sigma_j)\pm i\gamma_j\},\qquad
\sigma_j=\half+b_j>\half,\ \ \gamma_j>0 .
\]
The localized Weil form $Q$ --- the quadratic form whose positivity is $\RH$ --- is indefinite,
and its negative (Kre\u\i n) index $\kappa$ is the dimension of its maximal negative subspace
$\cN$; one has $\RH\iff\kappa=0$. The companion functional-analytic study of $Q$
\cite{Index} localizes $\kappa$ in the off-axis evaluation functionals attached to off-line
zeros and identifies $\kappa$ with the number of zeros of the structure function
$E_\xi(z)=\xi(\half+iz)$ in one open half-plane. The present paper takes that index as given
and develops its \emph{arithmetic} content.

The standing hypothesis throughout the conditional parts is:

\begin{quote}
\textbf{Hypothesis D.}\ $0<\kappa=2m<\infty$: there are finitely many off-line zeros, forming
$m\ge1$ Klein quartets.
\end{quote}

Hypothesis~D is strictly weaker than $\neg\RH$: it is the negation of $\RH$ \emph{together
with} finiteness of the off-line set. Finiteness is itself unconditionally open and weaker
than $\RH$ (it is not delivered by the positive-proportion or zero-density theorems
\cite{Index}); we assume it only to make the defect a finite, fully explicit object, and we
are careful to mark every statement that uses it.

The paper has three movements. \S\ref{sec:geometry} proves the geometry of $\cN$: parity
$\kappa=2m$, angular separation between orbits, and the Lorentz-plane decomposition.
\S\ref{sec:mef}--\S\ref{sec:wall} derive, under Hypothesis~D, the factorization
$\xi=\xi_0P_{4m}$ and the exact modified explicit formula $\psi=\psi_0-D_m$, analyze the
defect $D_m$, and prove the arithmetic-analytic gap that the companion $\zeta_0$ has no Euler
product --- the wall (RSRP). \S\ref{sec:twoscale}--\S\ref{sec:barrier} exhibit the two-scale
signature of an orbit, prove its compatibility with the unconditional prime number theorem and
with multiplicative chaos in the polynomially-bounded regime, and state the barrier theorem
isolating the three inputs that a proof of the exclusion principle $\kappa\in\{0,\infty\}$
would require. None of this assumes or proves $\RH$.

\section{The negative subspace: parity and Lorentz geometry}\label{sec:geometry}

We work in the reproducing-kernel Kre\u\i n space $\mathcal H$ of test functions on which $Q$
acts, with negative subspace $\cN=\{v:Q(v,v)<0\}$, $\dim\cN=\kappa$ \cite{Index}. For an
off-line orbit $\mathcal O_j$ write $\rho_j^+=\sigma_j+i\gamma_j$, $\rho_j^-=(1-\sigma_j)+
i\gamma_j$, and split the orbit by height into the upper sub-orbit $\mathcal O_j^+=\{\rho_j^+,
1-\bar\rho_j^+\}$ (both at height $+\gamma_j$) and the lower sub-orbit $\mathcal O_j^-=\{
\bar\rho_j^+,1-\rho_j^+\}$ (both at height $-\gamma_j$).

\begin{proposition}[The form restricted to one sub-orbit]\label{prop:restricted}
For test functions whose transform is concentrated near $+\gamma_j$, the contribution of
$\mathcal O_j^+$ to the Weil form is
\[
Q|_{\mathcal O_j^+}(f,f)=4\big[(\Re\fhat(\rho_j^+))^2-(\Im\fhat(\rho_j^+))^2\big],
\]
a two-dimensional form of signature $(1,1)$ whose unique negative direction $v_j^+$ is
characterized by $\Re\fhat(\rho_j^+)=0$. The lower sub-orbit $\mathcal O_j^-$ gives the
analogous form near $-\gamma_j$, with unique negative direction $v_j^-$.
\end{proposition}
\begin{proof}
The contribution of $\mathcal O_j^+=\{\rho_j^+,1-\bar\rho_j^+\}$ to $\sum_\rho\fhat(\rho)^2$ is
$\fhat(\rho_j^+)^2+\fhat(1-\bar\rho_j^+)^2$. The functional-equation/reality symmetries give
$\fhat(1-\bar s)=\overline{\fhat(s)}$, so $\fhat(1-\bar\rho_j^+)=\overline{\fhat(\rho_j^+)}$;
writing $\fhat(\rho_j^+)=u+iv$, the sum is $(u+iv)^2+(u-iv)^2=2(u^2-v^2)$, i.e.\ $4(u^2-v^2)$
after the archimedean normalization. The lower sub-orbit follows by conjugation.
\end{proof}

\begin{proposition}[Sub-orbit and inter-orbit orthogonality]\label{prop:orth}
Let $f^\pm$ have transforms supported in a $\delta$-neighborhood of $\pm\gamma_j$ in the
imaginary direction, $\delta<\gamma_j$. Then $Q(f^+,f^-)=O(e^{-c\gamma_j^2/\delta^2})$. More
generally, negative directions localized at distinct heights $\gamma_j\ne\gamma_k$ satisfy
$Q(v_j,v_k)=O(e^{-c(\gamma_j-\gamma_k)^2/\delta^2})$.
\end{proposition}
\begin{proof}
$Q(f^+,f^-)=\sum_\rho\fhat^+(\rho)\fhat^-(\rho)-B_{\mathrm{arith}}(f^+,f^-)$. The archimedean
cross term involves $\int\fhat^+\fhat^-\,\Omega$ with $\Omega$ smooth; the supports are
disjoint, so it vanishes. In the zero sum, a zero at height $\sim+\gamma_j$ contributes
$|\fhat^-(\rho)|=O(e^{-c\gamma_j^2/\delta^2})$ by concentration, and symmetrically; summing
over zeros (Riemann--von Mangoldt density) gives the bound. The inter-orbit estimate is
identical with $\gamma_j-\gamma_k$ in place of $\gamma_j$.
\end{proof}

\begin{theorem}[Parity of the index]\label{thm:parity}
$\kappa\equiv0\pmod2$; equivalently $\kappa=2m$, with $m$ the number of off-line Klein
quartets. If $\RH$ fails, then $\kappa\ge2$: the possible values of $\kappa$ are
$\{0,2,4,\dots\}$, and no odd value is consistent with the symmetry of $\zeta$. The same holds
for any primitive real Dirichlet $L$-function. With multiplicities, a zero of order $\ell$
contributes $2\ell$ to $\kappa$.
\end{theorem}
\begin{proof}
Each orbit $\mathcal O_j$ contributes, by Proposition~\ref{prop:restricted}, one negative
direction $v_j^+$ from $\mathcal O_j^+$ and one $v_j^-$ from $\mathcal O_j^-$. These are
independent: $Q(v_j^+,v_j^-)=O(e^{-c\gamma_j^2/\delta^2})$ by Proposition~\ref{prop:orth} with
$\delta<\gamma_j$, strictly smaller than $\sqrt{Q(v_j^+,v_j^+)\,Q(v_j^-,v_j^-)}$, so the two
are not proportional. The archimedean term and prime sum are non-negative, so every negative
direction arises from an off-line evaluation and is accounted for by this orbit structure;
hence $\dim\cN=2m$. This agrees with the structure-function count: the negative index equals
the number of zeros of $E_\xi$ in one open half-plane, which is exactly the two members of
each quartet on a fixed side, i.e.\ $2m$ \cite{Index}.
\end{proof}

\begin{remark}[Why two, and where the count is robust]\label{rem:two}
The two negative directions of an orbit live at the two heights $\pm\gamma_j$ and are
independent in the full (complex) Kre\u\i n space $\mathcal H$; on a space of \emph{even real}
test functions the reality $\fhat(\bar z)=\overline{\fhat(z)}$ identifies the two heights and a
quartet manifests there as a single Lorentzian block. The robust, representation-independent
count is the structure-function one, $\kappa=\#\{$zeros of $E_\xi$ in one half-plane$\}=2m$
\cite{Index}; the height decomposition above is the geometric realization of those two zeros.
\end{remark}

\begin{theorem}[Lorentz-plane decomposition]\label{thm:lorentz}
Under Hypothesis~D the negative subspace decomposes as an approximate orthogonal direct sum of
$m$ Lorentz planes,
\[
\cN\;\approx\;\bigoplus_{j=1}^m \mathbb R^{1,1}_j,
\]
where $\mathbb R^{1,1}_j=\operatorname{span}\{w_j^+,v_j^+\}$ carries signature $(1,1)$
($w_j^+$ the positive, $v_j^+$ the negative direction of $\mathcal O_j^+$), and the planes at
distinct heights are exponentially nearly orthogonal in $Q$ (Proposition~\ref{prop:orth}). The
approximation is exact in the limit of well-separated heights.
\end{theorem}
\begin{proof}
Each orbit contributes the signature-$(1,1)$ plane of Proposition~\ref{prop:restricted} at
height $+\gamma_j$ (and its mirror at $-\gamma_j$); by Proposition~\ref{prop:orth} the Gram
matrix between planes at $\gamma_j\ne\gamma_k$ is $O(e^{-c(\gamma_j-\gamma_k)^2/\delta^2})$,
so the total form is block-diagonal up to exponentially small couplings.
\end{proof}

\begin{corollary}[Parity-gap trichotomy]\label{cor:trichotomy}
The index satisfies $\kappa\in\{0\}\cup2\mathbb Z_{>0}\cup\{\infty\}$. Combined with the
exclusion of the finite non-zero (Pontryagin) regime by the finite-defect analysis
\cite{Index}, the operative dichotomy is between $\kappa=0$ ($\RH$) and the
arithmetically-loaded regimes $2\mathbb Z_{>0}\cup\{\infty\}$ studied below. We do not assert
$\kappa\ne$ any particular value; Theorem~\ref{thm:parity} only fixes the parity.
\end{corollary}

\section{The modified explicit formula}\label{sec:mef}

We now assume Hypothesis~D and read off the arithmetic consequences. The finiteness of $\Zoff$
makes the off-line part of the zero sum a genuine finite object.

\begin{proposition}[Factorization]\label{prop:factor}
Under Hypothesis~D,
\[
\xi(s)=\xi_0(s)\,P_{4m}(s),\qquad P_{4m}(s)=\prod_{\rho\in\Zoff}(s-\rho),
\]
where $P_{4m}$ is a degree-$4m$ polynomial with conjugate-symmetric roots and $\xi_0$ is
entire, shares the functional equation $\xi_0(s)=\xi_0(1-s)$, and has all its zeros on
$\Re s=\half$.
\end{proposition}
\begin{proof}
Divide the Hadamard product of $\xi$ by the finite product over $\Zoff$; the quotient $\xi_0$
is entire of order one with zero set $\Zcl$. The set $\Zoff$ is $V$-invariant, so $P_{4m}$ is
invariant under $s\mapsto1-s$ and under conjugation, whence $\xi_0=\xi/P_{4m}$ inherits the
functional equation and reality.
\end{proof}

\begin{definition}[Arithmetic defect and virtual Chebyshev function]\label{def:defect}
Under Hypothesis~D set
\[
D_m(x):=\sum_{\rho\in\Zoff}\frac{x^\rho}{\rho},\qquad
\psi_0(x):=x-\sum_{\rho\in\Zcl}\frac{x^\rho}{\rho}-\log(2\pi)-\half\log(1-x^{-2}).
\]
$D_m$ is real-valued ($\Zoff$ is conjugation-symmetric) and is a finite sum of $4m$ terms.
\end{definition}

\begin{theorem}[Modified explicit formula]\label{thm:mef}
Under Hypothesis~D, for every non-integer $x>1$,
\[
\psi(x)=\psi_0(x)-D_m(x),
\]
an exact identity. Writing $E(x)=\psi(x)-x$ and $E_0(x)=\psi_0(x)-x$, one has
$E(x)=E_0(x)-D_m(x)$ with $E_0(x)=O(x^{1/2+\varepsilon})$ automatically (all zeros of $\xi_0$
are on the line); hence the dominant growth of $E$ is carried entirely by the defect $D_m$.
\end{theorem}
\begin{proof}
Insert $\sum_\rho=\sum_{\Zcl}+\sum_{\Zoff}$ into the classical explicit formula
$\psi(x)=x-\sum_\rho x^\rho/\rho-\log(2\pi)-\half\log(1-x^{-2})$. Under Hypothesis~D the
off-line sum is finite, needing no regularization; rearranging gives the identity. The bound
on $E_0$ is the explicit formula for $\xi_0$, whose zeros all lie on the line.
\end{proof}

\section{Anatomy of the defect}\label{sec:anatomy}

Fix an orbit $\mathcal O_j$ and write $\theta_j^\pm=\arg\rho_j^\pm$.

\begin{proposition}[Single-orbit defect]\label{prop:Dj}
\[
D_j(x)=\underbrace{\frac{2x^{1/2+b_j}}{|\rho_j^+|}\cos(\gamma_j\log x-\theta_j^+)}_{\text{dominant}}
+\underbrace{\frac{2x^{1/2-b_j}}{|\rho_j^-|}\cos(\gamma_j\log x-\theta_j^-)}_{\text{subdominant}} .
\]
Consequently $D_m(x)=\sum_{j=1}^m\frac{2x^{1/2+b_j}}{|\rho_j^+|}\cos(\gamma_j\log x-\theta_j^+)
+O(x^{1/2-b_{\min}})$ as $x\to\infty$.
\end{proposition}
\begin{proof}
Pair conjugate zeros: $\frac{x^{\sigma_j+i\gamma_j}}{\sigma_j+i\gamma_j}+\frac{x^{\sigma_j-
i\gamma_j}}{\sigma_j-i\gamma_j}=\frac{2x^{\sigma_j}}{|\rho_j^+|}\cos(\gamma_j\log x-\theta_j^+)$,
and similarly for the $(1-\sigma_j)$ pair with exponent $1-\sigma_j=\half-b_j$. Collecting the
subdominant terms over all orbits gives the asymptotic.
\end{proof}

\begin{proposition}[Defect energy]\label{prop:energy}
For the logarithmic energy $\mathcal E_j(X)=\int_X^{2X}|D_j(t)|^2t^{-2}\,dt$,
\[
\mathcal E_j(X)\sim\frac{2^{2b_j}-1}{|\rho_j^+|^2\,b_j}\,X^{2b_j}\to\infty\quad(b_j>0),
\qquad \sum_{j=1}^m\mathcal E_j(X)\sim C_m\,X^{2b_{\max}} .
\]
\end{proposition}
\begin{proof}
The dominant part of $|D_j(t)|^2$ is $\frac{4t^{1+2b_j}}{|\rho_j^+|^2}\cos^2(\gamma_j\log t-
\theta_j^+)+O(t)$; the cosine-square averages to $\half$ (Riemann--Lebesgue), and
$\int_X^{2X}t^{2b_j-1}\,dt=\frac{(2^{2b_j}-1)}{2b_j}X^{2b_j}$. The total is dominated by the
largest $b_j$.
\end{proof}

\section{The arithmetic-analytic gap and the wall (RSRP)}\label{sec:wall}

The factorization $\xi=\xi_0P_{4m}$ is purely \emph{analytic}. We now show it is incompatible
with the \emph{multiplicative} structure of $\zeta$. Set
$\zeta_0(s)=\xi_0(s)/(\pi^{-s/2}\Gamma(s/2))$, the companion Dirichlet-type function; from
$\xi'/\xi=\xi_0'/\xi_0+P_{4m}'/P_{4m}$ one gets
\[
-\frac{\zeta_0'(s)}{\zeta_0(s)}=-\frac{\zeta'(s)}{\zeta(s)}+\sum_{\rho\in\Zoff}\frac{1}{s-\rho}.
\]

\begin{lemma}[Arithmetic incapacity of $\zeta_0$]\label{lem:noEuler}
Under Hypothesis~D:
\begin{enumerate}
\item $\zeta_0$ has a pole at each $\rho\in\Zoff$, hence poles inside the strip $0<\Re s<1$;
\item $\zeta_0$ is not representable as a Dirichlet series $\sum_n a_n n^{-s}$ with all
$a_n\ge0$;
\item $\zeta_0$ admits no Euler product $\prod_p(1-a_pp^{-s})^{-1}$ with $|a_p|\le1$;
\item the formal $\Lambda_0$ defined by $-\zeta_0'/\zeta_0=\sum_n\Lambda_0(n)n^{-s}$ is not
the von Mangoldt function of any standard arithmetic object.
\end{enumerate}
\end{lemma}
\begin{proof}
(1) The poles of $\zeta_0$ in the strip are the zeros of $P_{4m}$, i.e.\ $\Zoff$. (2) By
Landau's theorem, a Dirichlet series with non-negative coefficients that converges somewhere
has its abscissa of convergence as a real singularity and is otherwise holomorphic to its
right --- in particular it has no poles in $0<\Re s<1$; (1) forbids this. (3) An Euler product
of the stated form is holomorphic and zero-free for $\Re s>1$ and its continuation has no poles
in $\Re s>0$; (1) again forbids it. (4) $\Lambda\ge0$ produces non-negative Dirichlet
coefficients, excluded by (2).
\end{proof}

\begin{remark}
Lemma~\ref{lem:noEuler} is not a contradiction with Hypothesis~D. It says the Kre\u\i
n--Langer-type factorization is available in the analytic category (entire functions with a
functional equation) but incompatible with the arithmetic category (multiplicative
coefficients, Euler product) of $\zeta$. The wall below is the assertion that this
incompatibility is ultimately fatal to the existence of a finite non-zero index.
\end{remark}

\begin{definition}[The wall (RSRP)]\label{def:wall}
\emph{Resonant Spectral Rigidity of the Euler Product} (RSRP) is the statement: no entire
function $f$ satisfying the functional equation of $\xi$, carrying an Euler product
$f(s)=\tfrac{s(s-1)}2\pi^{-s/2}\Gamma(s/2)\prod_p(1-a_pp^{-s})^{-1}$ with $|a_p|\le1$, can have
$0<\kappa(f)<\infty$.
\end{definition}

\begin{theorem}[(RSRP) is a genuine, strictly weaker wall]\label{thm:wall}
\emph{(i)} $\RH\Rightarrow$ (RSRP) (vacuously, $\kappa=0$). \emph{(ii)} (RSRP)$\not\Rightarrow
\RH$ in isolation: it excludes only the finite non-zero index, leaving $\kappa=\infty$
untouched, so it is consistent with infinitely many off-line zeros. \emph{(iii)} (RSRP) is not
a consequence of the positivity wall (an analytic condition on $Q$), the de Branges/global
geometry wall (existence of a reproducing-kernel space), or the pair-correlation wall (a
statistical hypothesis on spacings): it is an \emph{arithmetic} incompatibility on the Euler
product, of a different type from all three.
\end{theorem}
\begin{proof}
(i) is immediate. (ii): (RSRP) forbids $0<\kappa<\infty$ but says nothing about $\kappa=\infty$,
which is consistent with $\neg\RH$; hence it cannot imply $\RH$. (iii): the three established
walls are, respectively, positivity of a quadratic form, existence of a geometric realization,
and a spacing statistic; (RSRP) is the non-existence of a meromorphic-with-Euler-product object
of finite index, an incompatibility between a polar divisor in the strip and multiplicativity.
These are logically independent formulations.
\end{proof}

\section{Two-scale structure of a single orbit}\label{sec:twoscale}

A single off-line orbit manifests at two distinct scales: an exponentially growing oscillation
in the Chebyshev variable $x$ (equivalently $t=\log x$), and a localized depression of
$\log|\zeta(\half+it)|$ near the ordinate $t=\gamma_j$.

\begin{theorem}[Hadamard depression]\label{thm:depression}
Under Hypothesis~D, the contribution of $\mathcal O_j$ to $\log|\zeta(\half+it)|$ is, up to
$O(1)$ archimedean constants,
\[
\delta_j\log|\zeta(\tfrac12+it)|=\log\!\big(b_j^2+(t-\gamma_j)^2\big)+\log\!\big(b_j^2+
(t+\gamma_j)^2\big)-\log|\rho_j^+|^2-\log|\rho_j^-|^2+O(1).
\]
In particular: at $t=\gamma_j$ the depth is $2\log b_j+O(1)$; the depression is localized with
half-width $b_j$; and for $|t-\gamma_j|\gg b_j$ it is $O(\log(|t-\gamma_j|/b_j))$.
\end{theorem}
\begin{proof}
The four Hadamard factors $1-s/\rho$ of $\mathcal O_j$ at $s=\half+it$ have moduli
$\sqrt{b_j^2+(t\mp\gamma_j)^2}/|\rho_j^\pm|$; taking logs of their product gives the displayed
formula (the convergence factors $e^{s/\rho}$ contribute $O(1)$). Setting $t=\gamma_j$ gives
depth $2\log b_j+O(1)$.
\end{proof}

\begin{proposition}[Average invisibility]\label{prop:cesaro}
If $\gamma_j\in[T,2T]$ then
\[
\frac1T\int_T^{2T}|\delta_j\log|\zeta(\half+it)||^2\,dt=O(b_j(\log b_j)^2/T)\to0 .
\]
The depression is invisible to the Cesàro average.
\end{proposition}
\begin{proof}
The integrand is $O((\log b_j)^2)$ on the window $[\gamma_j-b_j,\gamma_j+b_j]$ of width
$2b_j$, and $O((\log|t-\gamma_j|)^2)$ outside with convergent tail; dividing by $T$ gives the
estimate.
\end{proof}

\begin{theorem}[Structure-function bounds and their compatibility]\label{thm:struct}
Let $X(t)=(\psi(e^t)-e^t)e^{-t/2}$ and $S_2(\tau,T)=\frac1T\int_0^T|X(t+\tau)-X(t)|^2\,dt$.
\begin{enumerate}
\item \emph{(Lower bound, under Hypothesis~D.)} For $\tau\notin\frac{2\pi}{\gamma_j}\mathbb Z$,
\[
S_2(\tau,T)\ge\frac{4A_j^2\sin^2(\gamma_j\tau/2)}{2b_j}\cdot\frac{e^{2b_jT}-1}{T}+O(e^{b_jT}/
\sqrt T),\qquad A_j=\frac{2}{|\rho_j^+|},
\]
so $S_2(\tau,T)\gg e^{2b_jT}/T$.
\item \emph{(Upper bound, unconditional.)} By Korobov--Vinogradov, $S_2(\tau,T)\le C\,e^{T}
e^{-2c_1\sqrt T}$.
\item \emph{(Compatibility.)} The two bounds never contradict: $e^{2b_jT}/T\le Ce^Te^{-c_1
\sqrt T}$ holds for large $T$ because $2b_j-1<0$. A contradiction would require $b_j>\half$,
impossible since $\sigma_j<1$.
\end{enumerate}
\end{theorem}
\begin{proof}
(1) Under Hypothesis~D, $X=X_0-Y$ with $Y(t)=D_m(e^t)e^{-t/2}$ and, by the on-line bound
$\sum_{\Zcl}|\rho|^{-2}<\infty$, the $X_0$ structure function is $O(1)$; the dominant term of
$|Y(t+\tau)-Y(t)|^2$ from orbit $j$ is $4A_j^2e^{2b_jt}\sin^2(\gamma_j\tau/2)\cos^2(\cdots)$,
and averaging $\cos^2\to\half$ with $\int_0^Te^{2b_jt}dt\sim e^{2b_jT}/(2b_j)$ gives the bound.
(2) From $|\psi(x)-x|\le Cx\exp(-c_1\sqrt{\log x})$, $|X(t)|\le Ce^{t/2}\exp(-c_1\sqrt t)$, and
square-integrating. (3) is the displayed inequality.
\end{proof}

\section{Chaos compatibility and the barrier theorem}\label{sec:barrier}

\sloppy The $t$-scale depression must coexist with the value distribution of $\zeta$ on the line. The
Fyodorov--Hiary--Keating picture (established in part by Arguin--Bovier--Kistler, Najnudel,
Radziwi\l\l--\allowbreak Soundararajan) gives, for $t$ uniform in $[T,2T]$,
$\min_{[T,2T]}\log|\zeta(\half+it)|\sim-\log\log T+O(\log\log\log T)$.

\begin{theorem}[Chaos compatibility is regime-dependent]\label{thm:chaos}
With $T\asymp\gamma_j$ the depression depth $2|\log b_j|+O(1)$ compares with the chaos minimum
$\sim-c\log\gamma_j$ as follows:
\begin{enumerate}
\item if $b_j\ge\gamma_j^{-C}$ (polynomially bounded below), the depth is $O(\log\gamma_j)$
and is accommodated --- \emph{compatible};
\item if $b_j=e^{-\gamma_j^\alpha}$, $\alpha\in(0,1)$, the depth is $\sim2\gamma_j^\alpha\gg
\log\gamma_j$ --- \emph{potentially incompatible};
\item if $b_j\le e^{-c\gamma_j}$, the depth is $\sim2c\gamma_j$ --- \emph{incompatible} with
the chaos model.
\end{enumerate}
Since no unconditional lower bound on $b_j$ is known, uniform compatibility across all
$b_j\in(0,\half)$ cannot be asserted.
\end{theorem}
\begin{proof}
The depth is Theorem~\ref{thm:depression}; the minimum is the cited chaos result. In (1),
$O(C\log\gamma_j)$ matches $O(-\log\gamma_j)$; in (2)--(3), $\gamma_j^\alpha$ and $\gamma_j$
exceed $\log\gamma_j$ for large $\gamma_j$.
\end{proof}

\begin{remark}[A correction]\label{rem:correction}
Korobov--Vinogradov bounds $b_j$ from \emph{above} only ($b_j\le\half-c_0/(\log\gamma_j)^{2/3}$);
it gives no lower bound, hence no upper bound on $|\log b_j|$. Any claim that the depression
depth is $O(\log\log\gamma_j)$ is therefore unfounded: the correct conclusion is the
regime-dependent Theorem~\ref{thm:chaos}.
\end{remark}

\begin{theorem}[Barrier theorem]\label{thm:barrier}
Any proof of the exclusion principle $\kappa\in\{0\}\cup\{\infty\}$ from the present results
requires at least one of:
\begin{itemize}
\item[\emph{(F1)}] an unconditional bound $\psi(x)-x=O(x^{1/2+\varepsilon_0})$ for some fixed
$\varepsilon_0<\min_j b_j$, independent of Hypothesis~D;
\item[\emph{(F2)}] a property of the Euler product (valid for $\Re s>1$) incompatible with poles
of $-\zeta'/\zeta$ in $\Re s\in(\half,1)$ --- i.e.\ the wall (RSRP);
\item[\emph{(F3)}] a property of the multiplicative-chaos model incompatible with a depression
of $\log|\zeta(\half+it)|$ of the depth permitted by Theorem~\ref{thm:chaos}.
\end{itemize}
None of (F1), (F2), (F3) is currently available.
\end{theorem}
\begin{proof}
(F1) is necessary: without a better-than-Korobov--Vinogradov bound, the structure-function
estimates are compatible (Theorem~\ref{thm:struct}), so no contradiction arises from $D_m$.
(F2) is necessary: the Euler product for $\Re s>1$ gives no information on zeros in
$\Re s\in(\half,1)$, and the standard candidate invariants (Selberg second moment, Mertens
function, the Weil functional) are each either circular or reduce to a known wall. (F3) is
necessary by Theorem~\ref{thm:chaos}. Their status: (F1) is equivalent to controlling zero
locations, hence circular; (F2) is the wall (RSRP), unproven; (F3) is regime-dependent and
cannot be closed without a lower bound on $b_j$.
\end{proof}

\section{Conclusion}

Under the sole hypothesis that the off-line set is finite, a single Klein quartet of off-line
zeros is a fully explicit object with a fully explicit footprint: two negative directions in
the Weil form (Theorem~\ref{thm:parity}), a Lorentz-plane block in the negative subspace
(Theorem~\ref{thm:lorentz}), an exact defect term $D_m$ in the explicit formula
(Theorem~\ref{thm:mef}) growing like $x^{1/2+b_{\max}}$ with energy $X^{2b_{\max}}$
(Proposition~\ref{prop:energy}), and a logarithmic depression of $\log|\zeta(\half+it)|$ at
$t=\gamma$ of depth $2|\log b|$ (Theorem~\ref{thm:depression}). All of this is compatible with
the unconditional prime number theorem (Theorem~\ref{thm:struct}) and, in the
polynomially-bounded regime, with the value distribution of $\zeta$ on the line
(Theorem~\ref{thm:chaos}). The obstruction to excluding such a configuration is sharp and
named: the analytic factorization $\xi=\xi_0P_{4m}$ produces a companion with poles in the
strip and hence no Euler product (Lemma~\ref{lem:noEuler}), the incompatibility
crystallized as the wall (RSRP) (Theorem~\ref{thm:wall}); and any unconditional exclusion
must supply one of three precisely identified inputs (Theorem~\ref{thm:barrier}), none of which
is available. The contribution is a map of exactly where the finiteness of the index meets,
and resists, the arithmetic of $\zeta$. Nothing here assumes or proves the Riemann Hypothesis.

\begin{thebibliography}{99}
\bibitem{Index} D.~A.~Trejo Pizzo, \emph{The localized Weil form and its Kre\u\i n--Pontryagin
index: localization, obstruction, and rigidity}, preprint, 2026.
\bibitem{deBranges} L.~de Branges, \emph{Hilbert Spaces of Entire Functions}, Prentice--Hall,
Englewood Cliffs NJ, 1968.
\bibitem{Davenport} H.~Davenport, \emph{Multiplicative Number Theory}, 3rd ed., GTM \textbf{74},
Springer, 2000.
\bibitem{Landau} E.~Landau, \emph{\"Uber einen Satz von Tschebyschef}, Math.\ Ann.\ \textbf{61}
(1905), 527--550.
\bibitem{Ivic} A.~Ivi\'c, \emph{The Riemann Zeta-Function: Theory and Applications}, Dover, 2003.
\bibitem{FHK} Y.~V.~Fyodorov, G.~A.~Hiary, J.~P.~Keating, \emph{Freezing transition,
characteristic polynomials of random matrices, and the Riemann zeta function}, Phys.\ Rev.\
Lett.\ \textbf{108} (2012), 170601.
\bibitem{ABK} L.-P.~Arguin, D.~Belius, A.~J.~Harper, \emph{Maxima of a randomized Riemann zeta
function, and branching random walks}, Ann.\ Appl.\ Probab.\ \textbf{27} (2017), 178--215.
\end{thebibliography}

\end{document}
